\subsection{Benchmark and Primary Model}
\label{sec:methods_benchmark}

The benchmark is medical\_calculator\_eval (Hugging Face dataset
\textit{ekacare/medical\_calculator\_eval}, test split, $N=1066$
vignettes), released by EkaCare as part of the KARMA evaluation
framework~\citep{karma2024}. We use the public release without
modification.

Each item is a single-turn task: a case-narrative vignette in
natural-language prose, a target calculator, an \emph{expected
numerical value} for a designated primary output field, and a
per-item tolerance. There is no conversation --- the model receives
the vignette once and produces a single response. Each item also
carries a per-row \emph{confinement instruction} --- a prompt suffix
that instructs the model to reply with a single JSON object
containing the primary field --- which the adapter appends to the
vignette before generation. Scoring is binary: the harness extracts a
numerical prediction from the model's reply (preferring fenced JSON
blocks, falling back to the last number in the text), and the item is
marked correct iff
$|\text{predicted} - \text{expected}| \le \text{tolerance}$. Parse
failures and missing-field responses score as wrong. All accuracy
numbers in this paper are this strict tolerance-based binary
correctness, not partial-credit scores. Critically, the input is
free-text prose rather than structured fields --- the reason an LLM
extractor is required to populate the verifier's reference at all.
Structured-input benchmarks (e.g., EHR-keyed inputs) would remove
the extraction bottleneck this paper diagnoses. Curation details
(vignette authorship, validation, calculator coverage) are documented
in the KARMA release. The primary model under evaluation is
Qwen3-30B-A3B-Thinking-2507~\citep{qwen3_thinking}, served via
vLLM~\citep{vllm2023} on the text-completions endpoint with
temperature $0.0$, top-$p$ $1.0$, max generation length $4096$ tokens,
and a tool-calling loop bounded at $10$ turns. The primary model's
sampling is deterministic by construction (temperature zero, frozen
MCP schema dump). The extractor's sampling varies by condition and
is documented in \cref{sec:methods_prereg}.

The fact extractor used in all extractor-enabled conditions is Claude
Sonnet 4.6~\citep{claude_sonnet_4_6}, called with a single pinned
model snapshot for the duration of each run. The post-hoc verifier
used in exploratory condition D is also Sonnet 4.6. We chose Sonnet
rather than a sibling of Qwen3 to ensure the verifier's reference is
not produced by the same model family as the primary, which would
defeat the purpose of an external check.

\subsection{Compute and Serving Infrastructure}
\label{sec:methods_compute}

Qwen3 is served by a single vLLM process on an Azure NC40ads-H100-v5
instance (one NVIDIA H100 80\,GB GPU), running on Databricks Runtime
17.3 (GPU/ML build, Scala 2.13). Inference uses the text-completions
endpoint, not chat completions, so the custom tool-calling loop
driver can intercept generation at \verb|</tool_call>| stop tokens.
Per-condition runs are parallelized across vignettes via a process
pool. The worker count affects wall-clock only, not accuracy or token
counts --- each vignette is processed independently and the primary
model is deterministic at temperature zero --- so it is tuned for
throughput and recorded per run in the result-bundle manifest
(\cref{app:logging}). Reported latency comes from a separate
fixed-concurrency timing pass (\cref{sec:results_cost}), not from these
production runs. Sonnet calls are dispatched against
Anthropic's public API; no model weights other than Qwen3 run locally.
Per-condition GPU-hours and estimated cloud costs are recorded in
each result bundle's provenance file.

Because the worker count varies across conditions, per-vignette
wall-clock from the accuracy run is not comparable across conditions
(contention differs by load). Accuracy and token counts are
load-independent and are taken from the full run; \emph{latency} is
measured separately in a fixed-concurrency pass that re-runs a
$50$-vignette subset of every condition at a single worker, yielding
a consistent single-request (deployment-style) latency. All latency
numbers in \cref{sec:results_cost} come from that pass.

\subsection{Conditions: Design and Purpose}
\label{sec:methods_conditions}

The study is organized around a single experimental factor: the
\emph{architecture of the input-verification pipeline}. Nine
conditions are partitioned into three groups (\cref{tab:conditions}):
three \emph{anchors} that bound the room for verification to improve,
four \emph{primary study} conditions that vary the extractor
architecture under a fixed best-practice verifier, and two
\emph{exploratory} conditions that test alternative verifier
mechanisms (in-prompt self-verify and output-surface post-hoc check).

Each primary study condition is presented at its best-practice
configuration: schema-gated verifier comparison, query-style
intervention prompts, and prompt-only calibrated abstention. These
three hygiene properties are identical across all primary study
conditions and are documented in \cref{sec:methods_prereg}; they are
not factors under test. The naive un-gated and correction-style
variants used in earlier pilot runs are documented as design
motivation in \cref{sec:methods_pilots}.

\subsubsection{Capability Anchors (No Extractor)}
\label{sec:methods_anchors}

\textbf{A} disables tools entirely and uses no system prompt; it is
the capability floor, reporting what the primary model can answer
from internal knowledge alone.

\textbf{B} enables MCP tools with no system prompt; it is the
pre-registered baseline used for the apparatus gate
(\cref{sec:methods_gate}). If Sonnet's accuracy on B falls within
$\pm 3$ pp of the published 81.9\% reference, the harness is treated
as drift-free.

\textbf{B$'$} enables tools and adds a system prompt that instructs
the model to invoke a calculator rather than answering from prior
knowledge (verbatim prompt in \cref{lst:prompt_bprime}). B$'$
isolates the effect of forced tool use independent of any verifier
and serves as the comparator for the primary study conditions.

\subsubsection{Primary Study: Extractor Architecture}
\label{sec:methods_primary}

Four conditions share an identical verifier and intervention design;
they differ only in the extractor pipeline. All four use the B$'$
system prompt, the schema-gated deterministic verifier
(\cref{sec:methods_verifier}), query-style intervention prompts, and
prompt-only abstention.

\textbf{B$'$+E (upfront).} One Sonnet call per vignette, at run
start, populating the union of all calculator schemas
($\sim$500 fields). The extracted record is consulted for every
subsequent tool-call interception.

\textbf{B$'$+E (reactive).} One Sonnet call per
\toolname{medical\_calculator\_output} tool call, scoped to the
fields the active calculator declares. The extractor sees only the
field names the calculator requests and the case text; it never sees
the values the primary model proposed.

\textbf{B$'$+E (reactive + citations).} Reactive extraction with the
extractor returning, for each populated field, the exact source span
in the vignette that supports the value. The span must match a
substring of the vignette verbatim; if validation fails, the field
is treated as null (abstention).

\textbf{B$'$+E (reactive + k-shot).} Reactive extraction sampled $k=3$
times per tool call at temperature $0.7$ with three independent
seeds. The three samples are reduced per field by majority vote;
three-way disagreement yields null (abstention).

The architectural axis isolated by each pairwise contrast is given
in \cref{tab:conditions} and the contrast logic in
\cref{sec:methods_stats}.

\subsubsection{Exploratory: Other Verifier Mechanisms}
\label{sec:methods_exploratory}

Two conditions test verifier mechanisms outside the
extractor--verifier architecture family. Neither is in the primary
contrast family.

\textbf{C} (force-tool + prompt-only input self-verify) is the
in-prompt analog of the input verifier: it uses the same forced-tool
footing as the primary study (the B$'$ system prompt) and additionally
instructs the model to verify each input value against the case before
calling a tool (merged prompt in \cref{lst:prompt_c}). No external
extractor, no separate verifier --- the model checks its own inputs.
Running C on the B$'$ footing (rather than on plain tools) makes it a
matched comparator for the B$'$+E arms: the contrast
B$'$+E (reactive) vs C isolates the external extractor--verifier
against simply prompting the model to self-check, holding tool-forcing
constant. C is re-run under this definition; the earlier unforced
self-verify variant is superseded.

\textbf{D} (force-tool + post-hoc Sonnet output verifier) runs the
primary model on the same forced-tool footing as the input-verification
arms (the B$'$ system prompt), then sends \texttt{(case, candidate
answer)} to a Sonnet reviewer after the final answer is produced
(reviewer prompt in \cref{lst:prompt_d}). On flag, the model is
re-prompted once with the reviewer's one-sentence issue. D is the
closest analog to ``check the final answer'' verification approaches,
and running it on the B$'$ footing makes it a matched comparator for
the input-verification arms: B$'$+E checks the inputs \emph{before}
dispatch, D checks the answer \emph{after}, with tool-forcing held
constant. D is re-run under this definition; the earlier unforced
variant is superseded.

\begin{table}[!htbp]
  \centering
  \caption{Nine conditions. Primary study conditions share identical
  hygiene (schema-gated verifier, query intervention, prompt-only
  abstention); they vary only in the extractor pipeline.}
  \label{tab:conditions}
  \begin{tabular}{llllll}
    \toprule
    Condition & Group & Tools & Sys. prompt & Verifier & Extractor pipeline \\
    \midrule
    A                              & anchor       & no  & none       & none           & --                       \\
    B (baseline)                   & anchor       & yes & none       & none           & --                       \\
    B$'$                           & anchor       & yes & force-tool & none           & --                       \\
    B$'$+E (upfront)               & primary      & yes & force-tool & deterministic  & upfront full-schema       \\
    B$'$+E (reactive)              & primary      & yes & force-tool & deterministic  & reactive per-call         \\
    B$'$+E (reactive + citations)  & primary      & yes & force-tool & deterministic  & reactive + source spans   \\
    B$'$+E (reactive + k-shot)     & primary      & yes & force-tool & deterministic  & reactive + $k{=}3$ vote   \\
    C                              & exploratory  & yes & force-tool + self-verify & prompt-only    & --             \\
    D                              & exploratory  & yes & force-tool & post-hoc Sonnet& --                       \\
    \bottomrule
  \end{tabular}
\end{table}

\subsection{Pre-Registered Configuration}
\label{sec:methods_prereg}

The four primary study conditions are run under identical hygiene and
sampling configuration; the only inter-condition variation is the
extractor pipeline. All values below are pre-registered and locked
prior to the study run; no tuning to results is permitted.

\subsubsection{Baked-In Hygiene}

\textbf{Schema-gated verifier comparison.} The verifier compares
\toolname{arguments[``input\_data'']} against the extracted record
only on fields declared as \texttt{required} in the active
calculator's JSON Schema. Fields the extractor populated but the
calculator does not require are ignored at verify time. This
suppresses flags on fields that cannot affect the calculator's
output and was motivated by the pilot diagnostic in
\cref{sec:methods_pilots} showing 58.6\% of legacy-condition flags
fired on non-required fields.

\textbf{Query-style intervention.} On flag, the verifier injects a
prompt that presents the disagreement as a question to be resolved
against the vignette, not as an assertion of model error (verbatim
template in \cref{app:feedback_template}). The format names the
disagreeing field, the model's proposed value, and the extractor's
value, and asks the model to consult the case to determine which is
supported.

\textbf{Prompt-only calibrated abstention.} The extractor prompt
instructs Sonnet to return \texttt{null} for any field whose value
the case does not clearly state or imply. The verifier's existing
no-flag rule on missing extractor values
(\cref{sec:methods_verifier}) then applies. No confidence
thresholding; no post-hoc adjustment.

\subsubsection{Locked Run-Time Parameters}

\Cref{tab:prereg_params} fixes the parameters that govern sampling,
re-prompting, and failure handling. These are identical across all
primary study conditions; the only inter-condition variation lives
in the extractor pipeline row.

\begin{table}[!htbp]
  \centering
  \caption{Pre-registered run-time parameters for the primary study.
  Locked prior to the study run.}
  \label{tab:prereg_params}
  \begin{tabular}{ll}
    \toprule
    Parameter & Value \\
    \midrule
    Sonnet snapshot                  & pinned per-run; logged in result bundle \\
    Sonnet extraction temperature    & $0.7$ \\
    Sonnet extraction $k$ samples    & $1$ (primary 3 cond.); $3$ (k-shot cond.) \\
    Sonnet seed strategy             & fresh per call; no prompt caching \\
    Qwen3 temperature                & $0.0$ \\
    Tool-call loop turn cap          & $10$ \\
    Verifier re-prompt cap (total)   & $2$ per vignette \\
    Citation validity rule           & verbatim substring of vignette \\
    Citation validation failure      & treat field as null (abstain) \\
    k-shot voting rule               & per-field majority of 3 \\
    k-shot no-majority behavior      & treat field as null (abstain) \\
    Failure handling (incomplete vignettes) & drop from analysis; report dropped-$N$ per condition \\
    \bottomrule
  \end{tabular}
\end{table}

\subsubsection{Pilot-Batch Sanity Checks}

Before committing to the full $N=1066$ run, a 30-vignette pilot batch
verifies five properties of the implementation:
(1) k-shot samples diverge string-wise across the three seeds;
(2) the Qwen3 input for a given vignette in conditions B and B$'$ is
bit-identical (catches preprocessing drift);
(3) the re-prompt cap halts the loop when reached;
(4) the abstention prompt produces \texttt{null} on a vignette with
intentionally missing information;
(5) the citation validator rejects a paraphrased source span.
Failures on the pilot batch block the full run.

\subsection{Pilot Diagnostics (Design Motivation)}
\label{sec:methods_pilots}

The hygiene properties baked into the primary study were not
arbitrary defaults. Two pilot variants of B$'$+E (un-gated comparison
surface, correction-style intervention) ran prior to the primary
study and produced two findings that motivated the design.

First, under the un-gated upfront variant, the verifier compared
every field the extractor populated against the model's tool call,
regardless of whether the active calculator required the field.
Of $1{,}062$ total violations recorded across $N=1066$, $622$
($58.6\%$) fired on fields the active calculator's JSON Schema does
not declare as required; $181$ of those $622$ were on the
\verb|sex| field alone. The schema-gated comparison rule
(\cref{sec:methods_prereg}) eliminates this class of flag by
construction.

Second, the correction-style intervention prompt
(``Your planned tool call has an input that does not match the
patient case; re-read the case and call the tool again with
corrected inputs''') re-prompted the model with an assertion that
the verifier was correct. The query-style template
(\cref{app:feedback_template}) presents the same information without
that assertion, which lets the model adjudicate against the vignette
rather than defer to the verifier's reference.

These pilot results are reported in the result bundle for
reproducibility but are \emph{not} conditions in the primary study.
The primary study runs the best-practice versions only.

\subsection{The Schema-Gated Deterministic Verifier}
\label{sec:methods_verifier}

A note on naming. We call the verifier \emph{deterministic}, not
\emph{semantic}. Its comparison logic is purely structural: exact
match on enums, numeric tolerance on numbers, canonical-form match
on booleans. It does not reason about paraphrase, implication,
equivalent encodings, or context. The schema gate is the only
non-trivial filter and is purely structural as well: it consults the
active calculator's JSON Schema and restricts comparison to the
fields the calculator declares as \texttt{required}.

\Cref{fig:mechanism} shows the data flow. For each
\verb|<tool_call>| emitted by Qwen3, the verifier
(a) parses \toolname{arguments[``input\_data'']} into (field, value)
pairs;
(b) restricts the comparison set to fields the active calculator's
schema declares as required (\emph{schema gate});
(c) for each remaining pair, looks up the extractor's value for that
field;
(d) if the extractor has no value (abstention or missing),
\emph{skips} the field (no flag);
(e) if both have values, applies the deterministic comparison and
flags on mismatch.

% Apparatus schematic for §3.3.
%
% Routing rules:
%   - Main flow: straight horizontal/vertical only.
%   - Tool response: ONE right angle. UP from MCP, then LEFT into primary.east.
%     Use TikZ |- operator so the right angle is geometrically guaranteed
%     (no diagonal segments due to mismatched y-coordinates).
%   - Feedback wrap: leaves the BOTTOM of verify, wraps around the LEFT
%     of the diagram, arrives at the TOP of primary. All four corners
%     use a uniform clearance \pad so the halo around the diagram is
%     visually balanced.
%
% Label rules: every arrow gets a consistent verb-form gerund
%   (consults / dispatches / flags / tool response).
\begin{figure}[H]
\centering
\begin{tikzpicture}[
  font=\sffamily\small,
  >={Stealth[length=2.4mm]},
  every node/.style={align=center},
  box/.style={rectangle, rounded corners=2pt, draw=accentDark, line width=0.6pt,
              minimum height=11mm, minimum width=34mm, inner sep=4pt,
              fill=paleGray},
  modelbox/.style={box, fill=accentLight, draw=accent, line width=0.8pt},
  decisionbox/.style={diamond, draw=accent, line width=0.8pt, fill=white,
                      inner sep=2pt, aspect=2.4, minimum width=30mm,
                      minimum height=14mm},
  toolbox/.style={box, fill=white, draw=warmGray},
  arrow/.style={->, draw=accentDark, line width=0.6pt},
  feedback/.style={->, draw=condBprimeE, line width=0.7pt, dashed},
  cleanpath/.style={->, draw=condReactive, line width=0.7pt},
  ret/.style={->, draw=accentDark, line width=0.6pt},
]

% ── Grid coordinates ───────────────────────────────────────────────
\def\colL{0}        % left column x
\def\colM{5.6}      % middle column x
\def\colR{11.2}     % right column x
\def\rowT{4.5}      % top row y
\def\rowC{2.25}     % middle row y
\def\rowB{0}        % bottom row y

% Clearances for the feedback wrap-around halo.
% Three independent knobs so you can tune each side without affecting
% the others. Increase a value to push that side further from the
% diagram; decrease to bring it closer.
\def\padBelow{0.6}   % gap below verify before the path turns left
\def\padLeft{1.2}    % gap left of the leftmost column (this is the LEFT channel)
\def\padAbove{0.6}   % gap above primary before the path turns right and down

% ── Nodes ──────────────────────────────────────────────────────────
\node[box]      (case)      at (\colL,\rowT) {Patient vignette\\(natural language)};
\node[modelbox] (extractor) at (\colL,\rowC) {\textbf{Sonnet extractor}\\\footnotesize $M_e$: case $\to$ \texttt{\{field: value\}}};
\node[box]      (facts)     at (\colL,\rowB) {Patient facts\\(reference)};

\node[modelbox]    (primary) at (\colM,\rowT) {\textbf{Qwen3 primary}\\\footnotesize $\langle$tool\_call$\rangle$ generator};
\node[box]         (call)    at (\colM,\rowC) {Tool call payload\\\footnotesize\texttt{input\_data = \{...\}}};
\node[decisionbox] (verify)  at (\colM,\rowB) {\textbf{verify}\\\footnotesize field-by-field};

\node[toolbox] (mcp) at (\colR,\rowB) {MCP calculator\\(deterministic)};

% ── Main flow (all straight) ───────────────────────────────────────
\draw[arrow] (case)      -- (extractor);
\draw[arrow] (extractor) -- (facts);
\draw[arrow] (case)      -- (primary);
\draw[arrow] (primary)   -- (call);
\draw[arrow] (call)      -- (verify);
\draw[arrow] (facts)     -- node[above, font=\footnotesize\itshape,
                                  text=warmGray, yshift=0.5mm]
                            {consults} (verify);
\draw[cleanpath] (verify) -- node[above, font=\footnotesize\bfseries,
                                    text=condReactive] {dispatches}
                              (mcp);

% ── Tool response: ONE right angle via |- operator ─────────────────
% (mcp.north) |- (primary.east) automatically produces:
%   1. Vertical segment UP from mcp.north to primary.east.y
%   2. Right angle at the corner
%   3. Horizontal segment LEFT to primary.east
% No diagonals are possible because |- guarantees axis-aligned paths.
\draw[ret] (mcp.north) |- (primary.east);
% Label: anchored at center on the right vertical lane.
\node[font=\footnotesize\itshape, text=warmGray,
      anchor=center, rotate=90]
  at (\colR+0.30, \rowC) {tool response};

% ── Feedback wrap (left side, balanced halo with uniform \pad) ─────
% Path corners (all axis-aligned):
%   verify.south  → (verify.south.x, verify.south.y - \pad)   [DOWN]
%                 → (colL - boxhalf - \pad, same y)             [LEFT]
%                 → (same x, primary.north.y + \pad)            [UP]
%                 → (primary.north.x, same y)                   [RIGHT]
%                 → primary.north                              [DOWN]
\coordinate (verifyBelow) at ($(verify.south)+(0,-\padBelow)$);
\coordinate (leftLane)    at ($(case.west)+(-\padLeft,0)$);
\coordinate (topLane)     at ($(primary.north)+(0,\padAbove)$);

\draw[feedback]
  (verify.south)
    -- (verifyBelow)                           % DOWN
    -- (leftLane |- verifyBelow)               % LEFT to the left lane
    -- (leftLane |- topLane)                   % UP along the left lane
    -- (primary.north |- topLane)              % RIGHT to above primary
    -- (primary.north);                        % DOWN into primary.north

% Labels positioned along the path:
%   "flags" — short action label at the start of the dashed line
%             (just to the right of where it leaves verify.south).
%   "feedback re-prompt" — descriptive label on the left vertical lane,
%             mirror-image of "tool response" on the right lane.
\node[font=\footnotesize\bfseries, text=condBprimeE,
      anchor=west]
  at ($(verify.south)+(0.15,-0.3)$) {flags};
\node[font=\footnotesize\itshape, text=condBprimeE,
      anchor=center, rotate=90]
  at ($(leftLane)+(-0.30,-2.25)$) {feedback re-prompt};

\end{tikzpicture}
\caption{Apparatus overview. Right: a clean tool call is dispatched to
MCP. Dashed left: a verifier flag re-prompts the primary.}
\label{fig:mechanism}
\end{figure}


On flag, the verifier injects the query-style intervention template
(\cref{app:feedback_template}) into the model's prompt --- naming
the disagreeing fields, the value the model proposed, and the value
the extractor found --- and re-prompts. The MCP call is \emph{not}
dispatched. On clean, the call is dispatched and the tool response is
injected back into the prompt; the loop continues. On malformed JSON,
the verifier requests a correction without comparison. Re-prompts
per vignette are capped at the value in \cref{tab:prereg_params}.

Three properties of this design are critical. First, the verifier is
\emph{blind} to whether the extractor was right --- it only knows
whether the model and the extractor disagree on a required field both
populated; the extractor's value is treated as the reference by
construction. Second, the verifier's ``no value $\rightarrow$ no
flag'' rule conservatively suppresses false positives when the
extractor abstains. Third, the schema gate guarantees that flags on
fields the calculator's output does not depend on cannot occur, so
the verifier's harm cannot be attributed to spurious flags on
inactive fields.

\subsection{The Fact Extractor}
\label{sec:methods_extractor}

The extractor's vocabulary is generated deterministically from the
MCP per-calculator JSON Schemas via the schema dump. Because both
the extractor and the verifier consume the same schemas, the
verifier \emph{cannot ask about a field the extractor cannot in
principle emit}, and the extractor cannot emit a field the MCP
backend does not recognize. This vocabulary alignment is a
structural property of the design, not a runtime invariant we have
to police.

Four extractor pipelines are evaluated. All four share the same
prompt header (\cref{lst:prompt_extractor_header}), abstention rule,
and sampling temperature ($0.7$). They differ on three orthogonal
axes:

\begin{itemize}\setlength\itemsep{2pt}
  \item \textbf{Placement.} Where extraction happens relative to tool
        calls: upfront (full-schema, once per vignette) or reactive
        (scoped to the active calculator, once per tool call).
  \item \textbf{Output format.} Whether each extracted field is a
        bare value or a \texttt{(value, source\_span)} pair anchored
        in the vignette text.
  \item \textbf{Sampling.} Whether each field's value comes from a
        single sample ($k=1$) or a majority vote over $k=3$
        independent samples.
\end{itemize}

Each primary study condition is a single point in this space:

\begin{itemize}\setlength\itemsep{2pt}
  \item B$'$+E (upfront): upfront full-schema, bare values, $k=1$.
  \item B$'$+E (reactive): reactive scoped, bare values, $k=1$.
  \item B$'$+E (reactive + citations): reactive scoped, citation pairs, $k=1$.
  \item B$'$+E (reactive + k-shot): reactive scoped, bare values, $k=3$ vote.
\end{itemize}

The mechanism arms (citations, k-shot) anchor on reactive rather
than upfront for two reasons. First, $k=3$ over a $\sim$50K-token
upfront prompt is not a credible deployment configuration. Second,
each pairwise contrast against the reactive baseline isolates a
single axis change (\cref{sec:methods_stats}), preserving
orthogonality.

\subsubsection{Upfront Pipeline}
\label{sec:methods_extractor_upfront}

One Sonnet call per vignette, at run start, covering the full
$\sim$500-field union schema. The result is fixed for the duration of
the vignette and consulted on every tool-call interception. Across
all $\sim$70 calculators in MedAI MCP, the union of property names
exceeds 500 unique fields; the schema for any single calculator
typically declares 3--30. One representative calculator's JSON
Schema is reproduced verbatim in \cref{app:bmi_schema}.

\subsubsection{Reactive Pipeline}
\label{sec:methods_extractor_reactive}

One focused Sonnet call per
\toolname{medical\_calculator\_output} tool call. At each such call,
the adapter parses the keys of the model's proposed
\verb|input_data|, builds a focused Sonnet prompt restricted to those
fields, and runs extraction blind to the model's proposed values
(the extractor sees field \emph{names} and the case text, never the
values the model proposed). The result mutates the monitor's facts
for the duration of that tool call's verification. Per-vignette
cost scales with the number of \verb|_output| calls; each call is
small (typically 3--30 fields).

\subsubsection{Reactive + Citations Pipeline}
\label{sec:methods_extractor_citations}

Identical to the reactive pipeline, except each populated field is
returned as \texttt{\{value, source\_span\}} where
\texttt{source\_span} is a verbatim substring of the vignette text.
A post-extraction validator confirms the substring property; if
validation fails, the field is treated as null (abstention).
Validation rejects paraphrase, normalization, and reformatting; only
byte-identical spans are accepted. The verifier consults the
\texttt{value} field for comparison but logs both for analysis.

\subsubsection{Reactive + k-Shot Pipeline}
\label{sec:methods_extractor_kshot}

Identical to the reactive pipeline, except each focused Sonnet call
is dispatched $k=3$ times with three independent seeds and identical
prompt at temperature $0.7$. Samples are reduced per field by
majority vote: a value present in $\ge 2$ of $3$ samples wins;
three-way disagreement yields null (abstention). The three samples
are dispatched sequentially, so the per-call extraction latency of
this pipeline is approximately $3\times$ that of the single-sample
reactive pipeline; the cost--latency reporting
(\cref{sec:results_cost}) reflects this.

The only architectural differences across the four pipelines are
\emph{when} extraction happens, \emph{what} the extractor returns,
and \emph{how many independent samples} are taken. The verifier, the
monitor, the MCP transport, the loop driver, and the response schema
are byte-identical across all four.

\subsection{Pre-Registered Statistical Analysis}
\label{sec:methods_stats}

The primary contrast family targets the experimental factor
(extractor architecture). Six paired comparisons are pre-registered;
each isolates a single design axis:

\begin{enumerate}\setlength\itemsep{2pt}
  \item B vs A --- does tool access help?
  \item B$'$ vs B --- does forced tool use help, without verifier?
  \item B$'$+E (upfront) vs B$'$ --- does upfront verifier-guided
        extraction help over forced-tool baseline?
  \item B$'$+E (reactive) vs B$'$ --- does reactive verifier-guided
        extraction help over forced-tool baseline?
  \item B$'$+E (reactive + citations) vs B$'$+E (reactive) --- does
        citation grounding improve the reactive pipeline?
  \item B$'$+E (reactive + k-shot) vs B$'$+E (reactive) --- does
        $k$-shot voting improve the reactive pipeline?
\end{enumerate}

With family $\alpha = 0.05$ and Bonferroni correction over $n = 6$
tests, the per-test threshold is $\alpha = 0.00833$. The outcome is
paired binary correctness per vignette; the test is McNemar's exact
test~\citep{mcnemar1947} on the discordant pairs. Per-condition
accuracies are reported with Wilson~\citep{wilson1927} 95\%
confidence intervals. Effect sizes ($\Delta$ accuracy in percentage
points) are reported alongside $p$-values.

Four secondary contrasts are reported outside the primary family and
are not Bonferroni-corrected: (i) B$'$+E (upfront) vs B$'$+E
(reactive), isolating placement holding all other axes fixed;
(ii) C vs B$'$, testing whether prompting the model to self-check its
inputs helps over the no-check forced-tool baseline; (iii) B$'$+E
(reactive) vs C, testing whether the external extractor--verifier
beats prompt-only self-check on the matched forced-tool footing; and
(iv) D vs B$'$, the post-hoc output-check comparator on the matched
forced-tool footing. All four comparators (C, the B$'$+E arms, and D)
now sit on the same forced-tool footing as B$'$, so they can be read
as alternative verification methods against a common no-check baseline;
the legacy unforced C-vs-B and D-vs-B comparisons are superseded.

\subsection{Apparatus Gate}
\label{sec:methods_gate}

Before accepting primary results, we replicated Sonnet on Condition
B and required the resulting accuracy to fall within
$[0.789, 0.849]$ --- $\pm 3$ pp of EkaCare's published 81.9\%
reference for Claude Sonnet 4.6 with tool access on
medical\_calculator\_eval~\citep{ekacare_error_analysis}. The gate
threshold was specified in advance in the locked pre-registration
config (recorded in each result bundle's manifest, \cref{app:logging})
and is not adjustable post-hoc. Its purpose is drift
detection on the harness, dataset, prompt rendering, and MCP
transport. It is not a validation of the primary Qwen3 results ---
we assume the harness behaves identically for the primary model
conditional on the gate.
